\uuid{QbuB}
\titre{Limite}
\auteur{bodin ; hanani ; mzari}
\organisation{exo7}
\licence{CC BY-NC-SA 4.0}
\datecreate{2020-01-20}
\niveau{L1}
\module{Analyse}
\chapitre{Suite}
\sousChapitre{Autre}
\difficulte{2}
\theme{Suites, limite}
\exo7qcmid{lille-1.suites.011}

\contenu{
\begin{question}
\qclassification{121.00}
\qtags{motcle=limite}
Soit $\displaystyle u_n=\frac{2^{n+1}-3^{n+1}}{2^n+3^n}$ et $\displaystyle v_n=\frac{n2^{2n}-3^n}{n2^{2n}+3^n}$. Quelles sont les bonnes réponses ?
\begin{answers}
    \good{$\displaystyle \lim _{n\to +\infty}u_n=-3$ et $\displaystyle \lim _{n\to +\infty}v_n=1$}
    \bad{$\displaystyle \lim _{n\to +\infty}u_n=+\infty$ et $\displaystyle \lim _{n\to +\infty}v_n=+          \infty$}
    \bad{$\displaystyle \lim _{n\to +\infty}u_n=-3$ et $\displaystyle \lim _{n\to +\infty}v_n=+\infty$}
    \bad{$\displaystyle \lim _{n\to +\infty}u_n=-\infty$ et $\displaystyle \lim _{n\to +\infty}v_n=1$}
\end{answers}
\begin{explanations}
D'abord, $\displaystyle u_n=\frac{3^{n+1}\times \left[\left(\frac{2}{3}\right)^{n+1}-1\right]}{3^n\times \left[\left(\frac{2}{3}\right)^{n}+1\right]}=3\frac{\left(\frac{2}{3}\right)^{n+1}-1}{\left(\frac{2}{3}\right)^n+1}$. Or $\displaystyle \left(\frac{2}{3}\right)^n$ est le terme général d'une suite géométrique de limite $0$, donc $\displaystyle \lim _{n\to +\infty}u_n=3\frac{0-1}{0+1}=-3$. De même,
$$\displaystyle v_n=\frac{n(2^2)^n-3^n}{n(2^2)^n+3^n}=\frac{n4^n-3^n}{n4^n+3^n}=\frac{n4^n\times \left[1-\frac{1}{n}\left(\frac{3}{4}\right)^{n}\right]}{n4^n\times \left[1+\frac{1}{n}\left(\frac{3}{4}\right)^{n}\right]}=\frac{1-\frac{1}{n}\left(\frac{3}{4}\right)^{n}}{1+\frac{1}{n}\left(\frac{3}{4}\right)^{n}}.$$
Donc $\displaystyle \lim _{n\to +\infty}v_n=1$ car $\displaystyle \lim _{n\to +\infty}\frac{1}{n}=0$ et$\displaystyle \lim _{n\to +\infty}\left(\frac{3}{4}\right)^{n}=0$.
\end{explanations}
\end{question}
}
